\documentclass[12pt]{article}
\usepackage{fullpage}
\usepackage{enumitem}
\usepackage{graphicx}
% \usepackage{movie15}
\usepackage{media9}
\usepackage{listings}
\lstset{language=c, numbers=left}
\usepackage{addfont}
\addfont{OT1}{d7seg}{\dviiseg}
\addfont{OT1}{deseg}{\deseg}

\begin{document}
\title{Using C to Perform I/O on Simulated Hardware}
\author{}
\date{Part 1 due the week of November 9 at 11:59pm on your lab day\\
Parts 2 \& 3 due the week of November 16 at 11:59pm on your lab day}
\maketitle

Due to the circumstances of this semester, we will not be using Altera
single board computers for labs. Instead, we will use a simulated single
board computer (SSBC). \textbf{This lab is designed for use on the csce.unl.edu
Linux server.}

The accompanying document, \textit{ssbc.pdf} describes how to use the SSBC
library: how to interact with the simulated SBC and how to use its inputs and
outputs in a program.

The file archive \textit{lab6.tar} includes six files:
\begin{description}
    \item [ssbc.h] The header file for the SSBC library.
    \item [Makefile] The Makefile to build the programs for this lab. You can build the programs individually, or you can build all of them using the command \texttt{make all}.
    \item [demo.c] A program that demonstrates using the SSBC library. Run the
        command \texttt{make demo} to build the program, and then run it with
        the command \texttt{./demo}. This file is described line-by-line in
        \textit{ssbc.pdf}
    \item [warmup.c] Starter code for part 1 of this lab. Run the command
        \texttt{make warmup} to build the program, and then run it with the
        command \texttt{./warmup}.
    \item [poll-calculator.c] Starter code for part 2 of this lab. Run the
        command \texttt{make poll-calculator} to build the program, and then
        run it with the command \texttt{./poll-calculator}.
    \item [interrupt-calculator.c] Starter code for part 3 of this lab. Run the command\\ \texttt{make interrupt-calculator} to build the program, and then run it with the command \texttt{./interrupt-calculator}.
\end{description}

\textbf{\textit{Note: }Do NOT use \texttt{printf} to print debugging
information. Use \texttt{ssbc\_print} instead.}

\section{Familiarize Yourself with the Simulated Single Board Computer}

\textbf{20 points}

Edit \textit{warmup.c} to write a program that uses the SSBC's input and output
registers to meet the following specification:

\begin{itemize}
    \item When the left-most toggle switch is in the ``on'' position, the left
        three seven-segment displays show {\dviiseg 231}; when the
        left-most toggle switch is in the ``off'' position, the left three
        seven-segment displays are blank.
    \item When the right-most toggle switch is in the ``on'' position, the
        right-most seven-segment display displays the number corresponding to
        the last number button that was pressed. When the right-most toggle
        switch is in the ``off'' position, the right-most seven-segment display
        is blank. If a number button has not yet been pressed since launching
        the program, the right-most seven-segment display will be blank, even
        if the toggle switch is in the ``on'' position.
    \item Toggling the middle two toggle switches has no effect.
\end{itemize}

\textit{Hint:} You may want to create a ten-element array of 8-bit bit vectors
to map a decimal digit to the bit vector needed to display that digit on a
single seven-segment display.

\vspace{0.5cm}

Rubric:
\begin{description}
    \item[\underline{\hspace{1cm}} +4] {\dviiseg 231} displays on the three
        left displays when the left-most toggle switch is in the ``on''
        position.
    \item[\underline{\hspace{1cm}} +2] The left three displays are blank when
        the left-most toggle switch is in the ``off'' position.
    \item[\underline{\hspace{1cm}} +2] The right-most display is blank if
        a number button has not yet been pressed, regardless of the position of
        the right-most toggle switch.
    \item [\underline{\hspace{1cm}} +4] If a number button has been pressed,
        then the right-most display shows the number of the most-recent number
        button pressed when the right-most toggle switch is in the ``on''
        position.
    \item[\underline{\hspace{1cm}} +2] The right-most display is blank when
        the right-most toggle switch is in the ``off'' position.
    \item[\underline{\hspace{1cm}} +1] The behaviors for the left-most toggle
        switch and for the right-most toggle switch can happen at the same time
        (\textit{e.g.}, {\dviiseg 231} can be made to appear and disappear
        while also updating the number entered from the keypad).
    \item[\underline{\hspace{1cm}} +1] Nothing happens as a result of toggling
        either of the middle two toggle switches.
    \item[\underline{\hspace{1cm}} +2] The SSBC screen does not get corrupted
        under casual use of the program.
    \item[\underline{\hspace{1cm}} +2] The SSBC remains responsive to inputs
        throughout using the program.
\end{description}

\section{Use Polling to Detect Inputs} \label{polling}

\textbf{35 points}

It may not have been obvious while you were writing your warmup program, but
you can not tell when the user has pressed the same number button twice in a
row by simply reading the BCD value from the number pad's register. This is
obvious if you run the demonstration program and pay attention to the debug
window. In Part~\ref{polling}, you will overcome this problem by
\textit{polling} the number pad register's \textit{dirty bit}.

``Dirty bit'' is a bit used to indicate that something has changed. Here, the
dirty bit indicates that the user has pressed a number button. ``Polling'' is a
technique to periodically read a value to determine if it has changed.

\begin{figure}
    \centering
    \includegraphics[width=10cm]{AreWeThereYet}
    \caption{Polling. \tiny Image by 20th Century Fox Television}
\end{figure}

Edit \textit{poll-calculator.c} to write a program that uses the SSBC's input
and output registers to meet the following specification:

\begin{itemize}
    \item Functionality: implement a simple 2-function decimal integer
        calculator.
    \begin{itemize}
        \item Initially, $operand1$ and $operand2$ will hold the
            value $0$.
        \item When all toggle switches are in the ``off'' position, the
            seven-segment displays will be blank.
        \item When toggle switch 3 (the switch that is toggled by the `a' key)
            is moved to the ``on'' position, the program starts to build
            $operand1$. When toggle switch 3 is moved to the ``off''
            position, $operand1$ takes the value that was built.
        \item When toggle switch 2 (the switch that is toggled by the `s' key)
            is moved to the ``on'' position, the program starts to build
            $operand2$. When toggle switch 2 is moved to the ``off''
            position, $operand2$ takes the value that was built.
        \item When building an operand:
        \begin{itemize}
            \item The right-most seven-segment display (the ``ones'' place)
                will always display a digit. The remaining seven-segment
                displays will not display leading $0$s.
            \item The value being built will initially be $0$, and so {\dviiseg
                0} will be initially displayed.
            \item When the first number button is pressed after starting to
                build an operand, the ``ones'' place will display that digit.
                For example, if the user presses `5', then {\dviiseg 5} will be
                displayed.
            \item There will never be leading $0$s in the ``tens,''
                ``hundreds,'' and ``thousands'' places.
            \item When the second, third, and fourth number button is pressed
                (if a second, third, or fourth number button is pressed) then
                the existing digits will shift one decimal digit to the left,
                and the newly-entered number will go in the ``ones'' place. For
                example, if the user presses `5', `3', `0', and `9', then the
                display will show these values in sequence: \\
                \begin{tabular}{r}
                {\dviiseg 0} \\
                {\dviiseg 5} \\
                {\dviiseg 53} \\
                {\dviiseg 530} \\
                {\dviiseg 5309}
                \end{tabular}
            \item There is no way to build a negative operand.
            \item Except as noted in Part~\ref{bonus}, you may assume that
                operands are representable with four or fewer decimal digits.
                Unless you implement the bonus, we are not specifying the
                correct behavior if the user presses a fifth number button.
        \end{itemize}
        \item Number button presses will be ignored except when building
            operands.
        \item When toggle switch 1 (the switch that is toggled by the `d' key)
            is moved to the ``on'' position, the seven-segment displays will
            display the sum of $operand1 + operand2$. If the sum is $0$ then
            {\dviiseg 0} shall be displayed; otherwise, there shall be no
            leading $0$s. Except as noted in Part~\ref{bonus}, you may assume
            that the sum is representable with four or fewer decimal digits.
            Unless you implement the bonus, we are not specifying the correct
            behavior for sums greater than $9,999$.
        \item When toggle switch 0 (the switch that is toggled by the `f' key)
            is moved to the ``on'' position, the seven-segment displays will
            display the difference of $operand1 - operand2$. If the difference
            is $0$ then {\dviiseg 0} shall be displayed; otherwise, there shall
            be no leading $0$s. Except as noted in Part~\ref{bonus}, you may
            assume that the difference is non-negative; unless you implement
            the bonus, the correct behavior for negative differences is
            unspecified.
        \item You may assume that at most one toggle switch will ever be in the
            ``on'' position. We are not specifying the correct behavior if two
            toggle switches are in the ``on'' position.
    \end{itemize}
    \item Implementation
    \begin{itemize}
        \item Poll the number pad register's dirty bit to determine when a
            number button has been pressed. The SSBC library will place a $1$
            in the dirty bit whenver a number button has been pressed.
        \item Afer reading the number pad register's BCD value, clear the dirty
            bit by setting it to $0$.
        \item You may use any built-in C operations.
        \item \textit{Hint: } You may want to write a function that will
            convert a integer into the bit vector for the seven-segment
            displays' register. The array to generate bit vectors for a single
            digit from the warmup exercise will also be useful.
    \end{itemize}
\end{itemize}

Rubric:
\begin{description}
    \item[\underline{\hspace{1cm}} +1] The displays are blank when all toggle
        switches are in the ``off'' position
    \item[\underline{\hspace{1cm}} +1] $Operand1$ with with one digit can be
        built.
    \item[\underline{\hspace{1cm}} +2] $Operand1$ with non-repeating digits can
        be built, \textit{e.g.}, $231$.
    \item[\underline{\hspace{1cm}} +3] $Operand1$ with repeating digits can be
        built, \textit{e.g.}, $2331$.
    \item[\underline{\hspace{1cm}} +1] When building $operand1$, the
        ``ones'' place is always displayed
    \item[\underline{\hspace{1cm}} +2] When building $operand1$,
        there are no leading $0$s in the ``tens,'' ``hundreds,'' or
        ``thousands'' places.
    \item[\underline{\hspace{1cm}} +1] When toggle switch 3 is moved to the
        ``off'' position, $operand1$ takes on the value built.
    \item[\underline{\hspace{1cm}} +9] It is possible to build $operand2$ in
        accordance with specification.
    \item[\underline{\hspace{1cm}} +1] When toggle switch 2 is moved to the
        ``off'' position, $operand2$ takes on the value built.
    \item[\underline{\hspace{1cm}} +5] When toggle switch 1 is moved to the
        ``on'' position, the sum of $operand1 + operand2$ is displayed without
        leading $0$s.
    \item[\underline{\hspace{1cm}} +5] When toggle switch 0 is moved to the
        ``on'' position, the difference of $operand1 - operand2$ is displayed
        without leading $0$s.
    \item[\underline{\hspace{1cm}} +3] Number button presses are detected by
        polling the dirty bit.
    \item[\underline{\hspace{1cm}} +1] The dirty bit is reset after obtaining a
        the value of the number button pressed.
\end{description}

\section{Use Interrupts to Detect Inputs} \label{interrupts}

\textbf{20 points}

In a more complex program, polling inputs can be unwieldy. An alternative is
interrupt-driven programming, in which your program responds to hardware
interrupts rather than repeatedly checking whether an update is available. A
mix of polling and interrupts is common when only some inputs generate an
interrupt. In Part~\ref{interrupts}, your program will detect number button
presses by reacting to a simulated hardware interrupt.

The simulated SBC simulates hardware interrupts with operating system signals.
Specifically, if your program sets the ``IRQ enable'' bit, then the SSBC libary
will raise \texttt{SIGUSR1} signal whenever a number button is pressed. The
SSBC library initially sets this signal to be ignored; however, if you register
a signal handler with \texttt{sigset} then your program can react to the signal
as though it were a hardware interrupt.

\begin{figure}
    \centering
    \includegraphics[width=10cm]{MomMomMom}
    \caption{Interrupts. \tiny Image by 20th Century Fox Television}
\end{figure}

Edit \textit{interrupt-calculator.c} to write a program that uses the SSBC's
input and output registers to meet the following specification:

\begin{itemize}
    \item Functionality: implement a simple 2-function decimal integer
        calculator using the same specification as the calulator from
        Part~\ref{polling}.
    \item Implementation
    \begin{itemize}
        \item Write a signal handler to take the appropriate action when a
            number button is pressed, based on the program's current state.
        \item Register that signal handler to execute whenever the
            \texttt{SIGUSR1} signal is raised.
        \item Do \textit{not} poll the number keypad's register. (You will
            still need to poll the toggle switches' register.)
        \item You may use any built-in C operations.
        \item \textit{Hint: } Much of the code you wrote for Part~\ref{polling}
            can be reused, but the overall structure will differ.
    \end{itemize}
\end{itemize}

Rubric:
\begin{description}
    \item[\underline{\hspace{1cm}} +0.33] The displays are blank when all toggle
        switches are in the ``off'' position
    \item[\underline{\hspace{1cm}} +4.5] It is possible to build $operand1$ in
        accordance with specification.
    \item[\underline{\hspace{1cm}} +0.33] When toggle switch 3 is moved to the
        ``off'' position, $operand1$ takes on the value built.
    \item[\underline{\hspace{1cm}} +4.5] It is possible to build $operand2$ in
        accordance with specification.
    \item[\underline{\hspace{1cm}} +0.33] When toggle switch 2 is moved to the
        ``off'' position, $operand2$ takes on the value built.
    \item[\underline{\hspace{1cm}} +2.5] When toggle switch 1 is moved to the
        ``on'' position, the sum of $operand1 + operand2$ is displayed without
        leading $0$s.
    \item[\underline{\hspace{1cm}} +2.5] When toggle switch 0 is moved to the
        ``on'' position, the difference of $operand1 - operand2$ is displayed
        without leading $0$s.
    \item[\underline{\hspace{1cm}} +5] Number button presses are detected by
        handling the \texttt{SIGUSR1} signal.
\end{description}

\section{Extra Credit} \label{bonus}

\textbf{Up to 15 bonus points}

The calculator specification left a few edge cases unspecified. For bonus
credit, implement some or all of these specifications for the edge cases.

\textbf{Note: } to receive bonus credit, you only need to implement the edge
cases in one of the calculator implementations; you do not need to implement
the edge cases in both calculator implementations. \textit{If you implement the
edge cases, indicate so in the header comments of the calculator implementation
that you want evaluated for bonus credit. \textbf{We will not evaluate a
calculator implementation for bonus credit that does not state that it has the
edge cases implemented.}}

\begin{itemize}
    \item \textbf{Handle an attempt to build a too-great operand (+5)}
    \begin{itemize}
        \item When building an operand, if the user presses a fifth number
            button, the operand being built is invalid. The seven-segment
            displays shall display {\dviiseg err}.
        \item The only way to clear this error is to move the operand's toggle
            switch to the ``off'' position. When the user does so, the original
            operand's value will be preserved. For example:
        \begin{itemize}
            \item $operand1$ holds the value $53$.
            \item The user moves toggle switch 3 to the ``on'' position and
                presses `4', `5', `2', `3' which causes the display will show
                these values in sequence: \\
                \begin{tabular}{r}
                {\dviiseg 0} \\
                {\dviiseg 4} \\
                {\dviiseg 45} \\
                {\dviiseg 452} \\
                {\dviiseg 4523}
                \end{tabular}
            \item If the user were to move toggle switch 3 to the ``off''
                position, then $operand1$ would take the value $4,523$.
                Instead, the user presses `7' which causes the display to show
                {\dviiseg err}.
            \item The user now moves toggle switch 3 to the ``off'' position,
                and $operand1$ retains the value $53$.
        \end{itemize}
        \item This edge case must work for both operands to receive any credit
            for this edge case.
    \end{itemize}
    \item \textbf{Handle a too-great sum (+3)} If, when toggle switch 1 is
        moved to the ``on'' position, the sum is greater than $9,999$ then the
        seven-segment displays shall display {\dviiseg err}.
    \item \textbf{Handle negative differences (+5)} If, when toggle switch 0 is
        moved to the ``off'' position, the difference is between $-1$ and $-
        999$ (inclusive) then the difference shall be displayed, including the
        negative sign. As with non-negative values, there shall be no leading
        $0$s. For example, subtracting $53-102$ would cause the seven-segment
        displays to display {\dviiseg -49}.
    \item \textbf{Handle a too-low difference (+2)} If, when toggle switch 0 is
        moved to the ``off'' position, the difference is less than $-999$ then
        the seven-segment displays shall display {\dviiseg err}.
    \begin{itemize}
        \item If you do not implement the 5-point ``Handle negative
        differences'' bonus, then you can receive the 2-point ``Handle a too-
        low difference'' by displaying {\dviiseg err} for any value less than
        $0$.
    \end{itemize}
\end{itemize}

\section{Deliverables}

\begin{enumerate}
\item \textit{warmup.c}, due the week of November 9 at 11:59pm on the day of
    your scheduled lab section
\item \textit{poll-calculator.c}, due the week of November 16 at 11:59pm on the
    day of your scheduled lab section
\item \textit{interrupt-calculator.c}, due the week of November 16 at 11:59pm
    on the day of your scheduled lab section
\end{enumerate}

\end{document}
